%!TEX root = TheoreticalPhysicsC.tex

\section{Lecture: Temperature}
Consider two systems $1$ and $2$ in contact, so that energy, but no particles, can be exchanged. The total system is isolated. The total energy is $E$. The probability of finding system $1$ at energy $E_1$ is given by
\begin{equation}\label{eq:probenergy}
	w(E_1) \propto N_{E_1}^{(1)}N_{E-E_1}^{(2)}
\end{equation}
Here $N_{E}^{(1)}$ and $N_{E}^{(2)}$ are to be regarded as functions of $E_1$. 

The reason lies in the total system: by (\ref{eq:H-tensor}) its Hamiltonian is the sum $\widehat{H} = \widehat{H}_1\otimes \mathbb{I} + \mathbb{I}\otimes \widehat{H}_2$, and energies add:
\begin{equation}
	\widehat{H}\left(|E_j\rangle\otimes |E_k\rangle\right) = (E_j+E_k)\,|E_j\rangle\otimes |E_k\rangle.
\end{equation}
Only the sum $E_j+E_k$ is fixed; what we are after is how it is divided up. At fixed $E_1$ each subsystem is microcanonical, and the total ensemble is the product of the two,
\begin{equation}
	\rho = \rho_1\otimes \rho_2, \qquad p_{jk} = \frac{1}{N_{E_1}^{(1)}}\cdot\frac{1}{N_{E-E_1}^{(2)}},
\end{equation}
and counting the states yields $w(E_1)$.

The maximum of the probability is found with the help of the derivative:
\begin{equation}
	\frac{d}{dE_1} w(E_1) = 0.
\end{equation}
Inserting (\ref{eq:probenergy}) we obtain
\begin{equation}
	N_{E-E_1}^{(2)}\left(\frac{d}{dE_1}N_{E_1}^{(1)}\right) - N_{E_1}^{(1)}\left(\frac{d}{d(E-E_1)}N_{E-E_1}^{(2)}\right) = 0.
\end{equation}
Rearranged:
\begin{equation}
	\frac{1}{N_{E_1}^{(1)}}\frac{d}{dE_1}N_{E_1}^{(1)} =  \frac{1}{N_{E-E_1}^{(2)}}\frac{d}{d(E-E_1)}N_{E-E_1}^{(2)}.
\end{equation}
It is no accident that logarithmic derivatives appear here: under tensor products the entropy is additive,
\begin{equation}
	S(\rho_1\otimes \rho_2) = S(\rho_1)+S(\rho_2).
\end{equation}
With 
\begin{equation}
	S_1(E_1) = k\log(N_{E_1}^{(1)})\quad \text{and} \quad S_2(E-E_1) = k\log(N_{E-E_1}^{(2)})
\end{equation}
this relation will take the following form
\begin{equation}
	\frac{dS_1(E_1)}{dE_1} = \frac{dS_2(E_2)}{dE_2}\bigg|_{E_2 = E-E_1}.
\end{equation}
The \emph{temperature}\footnote{It is really the \emph{inverse} temperature that is being defined here; that is how things have gone in physics.} $T$ is therefore defined by 
\begin{equation}
	\frac{1}{T} = \frac{\partial S(E)}{\partial E}
\end{equation}
According to this definition thermal equilibrium between systems $1$ and $2$ turns out to mean that the temperatures are equal.

% fig_vl21_kontakt.tex — Zwei Systeme in thermischem Kontakt: wo die
% Gesamtentropie ihr Maximum hat, und wie scharf dieses Maximum wird.
%
% WHY THIS FIGURE.  The lecture derives dS_1/dE_1 = dS_2/dE_2 by
% differentiating w(E_1) and then *defines* 1/T by it.  Two things stay
% invisible in the prose.  First, the condition is a statement about
% SLOPES: the sum is maximal where the rising curve climbs exactly as
% fast as the falling one drops — and that place is in general NOT at
% E_1 = E/2.  Second, the whole argument only means something because
% the maximum of w is sharp; how sharp is a scaling statement
% (width ~ 1/sqrt(N)) that a sentence cannot show.
%
% GEOMETRY FROM THE FORMULA.  Nothing is sketched.  The lecture leaves
% N_E abstract, so the plate fixes the simplest counting with a closed
% form — N oscillators sharing E/eps quanta (Einstein solid), whose
% state count in the Stirling form is
%
%       S(N,E)/k = N f(E/(N eps)),   f(x) = (1+x)log(1+x) - x log x.
%
% System 1 gets N_1 = 1, system 2 gets N_2 = 3 (three times as large),
% and the total energy is E = 2 eps.  Then, with u = E_1/E,
%
%       S_1/k = f(2u),        S_2/k = 3 f(2(1-u)/3).
%
% Because f'(x) = log(1 + 1/x), the stationarity condition
% dS_1/dE_1 = dS_2/dE_2 reduces to E_1/N_1 = E_2/N_2, so the maximum
% sits exactly at
%
%       E_1^* = E N_1/(N_1+N_2) = E/4,
%
% and the common slope there is f'(1/2) = log 3 (in units 1/eps); in the
% plotted variable u the two tangents therefore have the exactly mirrored
% slopes +2log3 and -2log3.  Every coordinate in panel (a) — the peak
% position, the tangency points, the tangent slopes, the value of
% S_1+S_2 at E/2 — is that arithmetic, not a drawing.
%
% Panel (b) is panel (a) exponentiated.  Scaling both systems by the
% same factor lambda (N_i -> lambda N_i, E -> lambda E) leaves u^* and
% the shape function g(u) = S_tot(u)/k untouched, so
%
%       w(E_1) ∝ N^{(1)} N^{(2)} = e^{S_tot/k} = e^{lambda g(u)},
%
% which is *exactly* the w(E_1) of eq. (eq:probenergy), normalised to its
% own maximum — panel (a) is the lambda = 1 member of that family.  The
% width shrinks like 1/sqrt(lambda): with g''(u^*) = -64/9 the Gaussian
% width is 0.375/sqrt(lambda).
%
% NUMERICS.  Writing that exponent as lambda*(S_1+S_2-Smax) is fatal at
% lambda = 100: pgfmath's log carries an absolute error of order 1e-4,
% the subtraction of two numbers near 3.82 keeps it, and multiplying by
% 100 turns it into a visibly ragged crest that falls ~2% short of 1.
% So the difference is formed analytically first.  With d = 2u - 1/2
% (zero at the maximum) the linear terms cancel exactly and
%
%   g(u) - g(u^*) = (3/2+d)log(1+d/(3/2)) - (1/2+d)log(1+d/(1/2))
%                 + (9/2-d)log(1-d/(9/2)) - (3/2-d)log(1-d/(3/2)),
%
% every logarithm now of the form log(1+small).  This vanishes at d = 0
% by construction, so all three curves reach 1 exactly, and its second
% derivative there is -16/9 in d, i.e. g''(u^*) = -64/9 as above.
%
% DECISIONS.
%   * colour carries which system a curve belongs to: system 1 blue,
%     system 2 orange, the sum red.  A tangent is dashed in the colour
%     of its own curve; construction lines that belong to no system
%     (the verticals at E_1^* and E/2) are neutral grey;
%   * panel (b) is an ordered ramp of the *same* red as the sum in (a),
%     because w = exp(S_tot/k) is that curve and nothing else;
%   * E_1 < 0 and E_2 < 0 are unreachable and hatched, not filled — the
%     constraint E_1 + E_2 = E is what makes this a one-parameter
%     problem at all;
%   * the y axis of (a) carries no ticks: the counting behind the shape
%     is a stand-in, only the slopes and the position of the maximum
%     are the physics.  It is labelled S, not S/k, so that a tangent's
%     slope is the lecture's own 1/T = dS/dE rather than 1/(kT).  The
%     numbers behind the curve are S/k; with no ticks and k > 0 that is
%     the same curve up to a scale, and the scale is what the axis name
%     fixes;
%   * E/2 is marked in grey solely so that one can see the maximum is
%     not there — temperatures equalise, energies do not;
%   * every label sits on its own object; the panels were laid out by
%     first computing the gaps between the curves, so no annotation
%     crosses a curve.

\begin{center}
\begin{tikzpicture}

% ─────────────────────────────────────────────────────────────
%  Shared abbreviations.  f(x) = (1+x)ln(1+x) - x ln x, evaluated at
%  x = 2u (system 1) and x = 2(1-u)/3 (system 2).
% ─────────────────────────────────────────────────────────────
\def\Sone{((1+2*x)*ln(1+2*x)-2*x*ln(2*x))}
\def\Stwo{(3*((1+(2-2*x)/3)*ln(1+(2-2*x)/3)-((2-2*x)/3)*ln((2-2*x)/3)))}
\def\ustar{0.25}          % E_1^*/E = N_1/(N_1+N_2)
\def\Sonestar{0.9547713}  % f(1/2)
\def\Stwostar{2.8643138}  % 3 f(1/2)
\def\Smax{3.8190851}      % 4 f(1/2)
\def\slp{2.1972246}       % d(S_1/k)/du = 2 log 3
\def\Shalf{3.6355738}     % S_1+S_2 at E_1 = E/2

% g(u) - g(u^*), formed analytically so that no large exponent ever
% amplifies a cancellation (see NUMERICS above).  d = 2u - 1/2.
\def\dd{(2*x-0.5)}
\def\Gdiff{( (1.5+\dd)*ln(1+\dd/1.5) - (0.5+\dd)*ln(1+\dd/0.5)
             +(4.5-\dd)*ln(1-\dd/4.5) - (1.5-\dd)*ln(1-\dd/1.5) )}

% ─────────────────────────────────────────────────────────────
%  (a)  The two entropies and their sum against E_1
% ─────────────────────────────────────────────────────────────
\begin{axis}[
    name=ent,
    tpcpotential,
    width=9.0cm, height=6.6cm,
    enlargelimits=false,
    xmin=-0.20, xmax=1.18, ymin=-0.45, ymax=4.80,
    xtick={0.25,0.5,1}, xticklabels={$E_1^{*}$,$E/2$,$E$},
    ytick=\empty,
    xlabel={$E_1$}, ylabel={$S$},
    clip=false,
    domain=0.0015:0.9985, samples=241,
    title={(a) \tpcifenglish{The two entropies and their sum}%
                            {Die beiden Entropien und ihre Summe}},
]

  % Unreachable: E_1 < 0 on the left, E_2 = E - E_1 < 0 on the right
  \fill[forbidden] (axis cs:-0.20,0) rectangle (axis cs:0,4.60);
  \fill[forbidden] (axis cs:1,0)     rectangle (axis cs:1.12,4.60);
  \node[lbl, text=birrengray, font=\scriptsize, rotate=90]
    at (axis cs:-0.10,2.3) {$E_1<0$};
  \node[lbl, text=birrengray, font=\scriptsize, rotate=90]
    at (axis cs:1.06,2.3)  {$E_2<0$};

  % Reference verticals: equal energies (dotted) and the maximum (dashed)
  \draw[classicallim]  (axis cs:0.5,0)     -- (axis cs:0.5,\Shalf);
  \draw[levellinegrey] (axis cs:\ustar,0)  -- (axis cs:\ustar,\Smax);

  % The three curves
  \addplot[wavefn, cgblue]              {\Sone};
  \addplot[wavefn, birrenorange]        {\Stwo};
  \addplot[line width=1.3pt, birrenred] {\Sone+\Stwo};

  % Tangents at E_1^*: slopes +2log3, -2log3 and 0 — mirror images
  \addplot[cgblue, line width=0.7pt, densely dashed, domain=0.14:0.56,
           samples=2] {\Sonestar+\slp*(x-\ustar)};
  \addplot[birrenorange, line width=0.7pt, densely dashed,
           domain=0.14:0.56, samples=2] {\Stwostar-\slp*(x-\ustar)};
  \addplot[birrenred, line width=0.7pt, densely dashed, domain=0.14:0.56,
           samples=2] {\Smax};

  % Points of tangency, and where the naive guess E_1 = E/2 would sit
  \node[statepoint, inner sep=1.3pt] at (axis cs:\ustar,\Sonestar) {};
  \node[statepoint, fill=birrenorange, inner sep=1.3pt]
    at (axis cs:\ustar,\Stwostar) {};
  \node[statepoint, fill=birrenred, inner sep=1.3pt]
    at (axis cs:\ustar,\Smax) {};
  \node[statepoint, fill=birrengray, inner sep=1.1pt]
    at (axis cs:0.5,\Shalf) {};

  % Labels, each in the gap belonging to its own curve
  \node[lbl, text=birrengray, font=\scriptsize, anchor=north east]
    at (axis cs:0.99,4.55) {$N_2=3\,N_1$};
  \node[lbl, anchor=south] at (axis cs:\ustar,3.92) {$T_1=T_2$};
  \node[lbl, text=birrenred, anchor=south west]
    at (axis cs:0.64,3.50) {$S_1+S_2$};
  \node[lbl, text=birrenorange, anchor=south west]
    at (axis cs:0.52,2.42) {$S_2(E-E_1)$};
  % The slope labels start just right of E_1^*, where the slopes are
  % taken.  Their heights are set by the two things they must not touch:
  % the label of S_2 clears the falling S_2 curve at its right end, the
  % label of S_1 clears the rising +1/T_1 tangent at the same place.
  \node[lbl, text=birrenorange, font=\scriptsize, anchor=north west]
    at (axis cs:0.28,2.36) {$-1/T_2$};
  \node[lbl, text=cgblue, font=\scriptsize, anchor=south west]
    at (axis cs:0.28,1.42) {$+1/T_1$};
  \node[lbl, text=cgblue, anchor=north west]
    at (axis cs:0.52,1.15) {$S_1(E_1)$};

  \node[lbl, anchor=north east] at (axis cs:-0.02,-0.05) {$0$};
\end{axis}

% ─────────────────────────────────────────────────────────────
%  (b)  The same total entropy, exponentiated: w(E_1) and its width
% ─────────────────────────────────────────────────────────────
\begin{axis}[
    tpcplot,
    at={(ent.south east)}, anchor=south west, xshift=13mm,
    width=5.4cm, height=6.6cm,
    xmin=0, xmax=1, ymin=0, ymax=1.4,
    xtick={0,0.25,1}, xticklabels={$0$,$E_1^{*}$,$E$},
    ytick={0,0.5,1},
    xlabel={$E_1$}, ylabel={$w/w_{\max}$},
    every axis y label/.style={at={(axis description cs:-0.19,0.5)},
                               rotate=90, anchor=south, font=\small},
    domain=0.0015:0.9985, samples=1201,
    legend pos=north east,
    legend style={font=\scriptsize, draw=light-gray},
    title={(b) \tpcifenglish{How sharp the maximum is}%
                            {Wie scharf das Maximum wird}},
]
  \draw[levellinegrey] (axis cs:\ustar,0) -- (axis cs:\ustar,1.02);

  \addplot[birrenred!50, line width=1.0pt, line join=round]
    {exp(max(1*\Gdiff,-16))};
  \addlegendentry{$N$}
  \addplot[birrenred!72, line width=1.0pt, line join=round]
    {exp(max(10*\Gdiff,-16))};
  \addlegendentry{$10\,N$}
  \addplot[birrenred, line width=1.1pt, line join=round]
    {exp(max(100*\Gdiff,-16))};
  \addlegendentry{$100\,N$}

  \node[lbl, text=birrengray, font=\scriptsize, anchor=north east]
    at (axis cs:0.99,1.00) {\tpcifenglish{width}{Breite}
                            $\propto 1/\sqrt{N}$};
\end{axis}

\end{tikzpicture}
\end{center}


The figure shows the maximum as an equality of slopes: $S_1$ rises exactly as fast as $S_2$ falls. It does not lie at $E_1=E/2$, and it becomes extremely sharp as the system grows.

\begin{merksatz}
	In equilibrium the temperatures equalise, not the energies -- derived, not assumed.
\end{merksatz}
